\lyseis{Γ2.1}{Μέση τιμή και διάμεσος}
\sol{1} α) $120:24=5$. β) $20:8=2{,}5$. γ) $10:40=0{,}25$.
\sol{2} Άθροισμα $=$ πλήθος $\times$ μέση τιμή. α) $30\cdot4=120$. β) $6\cdot2{,}5=15$. γ) $20\cdot0{,}1=2$.
\sol{3} Πλήθος $=$ άθροισμα $:$ μέση τιμή. α) $48:4=12$. β) $30:2{,}5=12$. γ) $20:0{,}5=40$.
\sol{4} α) $\dfrac{4+6+8+10+12}{5}=\dfrac{40}{5}=8$. β) $\dfrac{-2+0+3+5}{4}=\dfrac{6}{4}=1{,}5$ $^\circ$C. γ) $\dfrac{28}{4}=7$.
\sol{5} α) Στη σειρά: $1,3,5,7,9$. Διάμεσος $5$. β) Στη σειρά: $4,6,8,12$. Διάμεσος $\dfrac{6+8}{2}=7$. γ) Στη σειρά: $3,3,3,7,9,10$. Διάμεσος $\dfrac{3+7}{2}=5$.
\sol{6} α) Λάθος (π.χ. για $1,2$ η μέση τιμή είναι $1{,}5$). β) Σωστό. γ) Σωστό. δ) Λάθος (π.χ. για $1,2$ βγαίνει $1{,}5$). ε) Λάθος: η μέση τιμή είναι πάντα ανάμεσα στη μικρότερη και τη μεγαλύτερη τιμή.
\sol{7} Γινόμενα: $0\cdot4=0$, $1\cdot6=6$, $2\cdot8=16$, $3\cdot2=6$. Άθροισμα γινομένων $28$, σύνολο παρατηρήσεων $4+6+8+2=20$. Μέση τιμή $28:20=1{,}4$.
\sol{8} Άθροισμα $6+7+5+8+9+7=42$, μέση τιμή $42:6=7$. Στη σειρά: $5,6,7,7,8,9$, διάμεσος $\dfrac{7+7}{2}=7$. Και τα δύο μέτρα βγαίνουν $7$.
\sol{9} α) $9\cdot12=108$ τριαντάφυλλα. β) Νέο άθροισμα $108+22=130$, νέο πλήθος $10$, μέση τιμή $130:10=13$.
\sol{10} Για μέσο όρο $17$ στα πέντε τεστ χρειάζεται άθροισμα $5\cdot17=85$. Έως τώρα $16+18+15+17=66$. Στο πέμπτο χρειάζεται $85-66=19$.
\sol{11} Σύνολο $24$ παιδιά. Γινόμενα: $0,4,14,21,16,5$ με άθροισμα $60$. Μέση τιμή $60:24=2{,}5$ βιβλία.
\sol{12} Κέντρα: $2{,}5$, $7{,}5$, $12{,}5$, $17{,}5$. Γινόμενα: $2{,}5\cdot3=7{,}5$, $7{,}5\cdot8=60$, $12{,}5\cdot6=75$, $17{,}5\cdot3=52{,}5$. Άθροισμα $195$, σύνολο $20$. Μέση τιμή $195:20=9{,}75$ ευρώ.
\sol{13} Το νέο άθροισμα είναι $50+10=60$ και το πλήθος $6$: η μέση τιμή μένει $10$ (προσθέσαμε τιμή ίση με τη μέση τιμή). Με $16$: $\dfrac{50+16}{6}=\dfrac{66}{6}=11$.
\sol{14} Παράδειγμα: $3,5,7,12,23$. Η μεσαία είναι το $7$ και το άθροισμα $3+5+7+12+23=50$, άρα η μέση τιμή είναι $10$. Δεν είναι μοναδική: υπάρχουν πολλές πεντάδες που ικανοποιούν τις συνθήκες (π.χ. $1,6,7,13,23$).
\sol{15} Όχι. Τα δύο τμήματα δεν έχουν τον ίδιο αριθμό μαθητών, άρα το Β «βαραίνει» περισσότερο. Άθροισμα βαθμών: $10\cdot14+20\cdot17=140+340=480$. Σύνολο μαθητών $30$. Σωστός μέσος όρος $480:30=16$.
\sol{16} Μέση τιμή: $\dfrac{900+1000+1100+1000+5000}{5}=\dfrac{9000}{5}=1800$ €. Στη σειρά: $900,1000,1000,1100,5000$, διάμεσος $1000$ €. Ο τυπικός μισθός είναι γύρω στα $1000$ €. Η μέση τιμή ($1800$) τραβήχτηκε από τον μισθό των $5000$ και δεν μοιάζει με τον μισθό κανενός από τους τέσσερις. Η διάμεσος περιγράφει καλύτερα.
\sol{17} Ελεύθερη εργασία. Για παράδειγμα: ώρες $8,9,8,7,9,10,9$ έχουν άθροισμα $60$ και μέση τιμή $60:7\approx8{,}6$. Στη σειρά: $7,8,8,9,9,9,10$, διάμεσος $9$.
\sol{18} Ελεύθερη εργασία. Ο πίνακας έχει τιμές (πλήθος κινητών) και συχνότητες που αθροίζουν στο $10$. Η μέση τιμή είναι: άθροισμα των γινομένων τιμή $\times$ συχνότητα, διά $10$.
